\input{../tex-inputs/latex-leseansicht-vorspann}

\section[Arthur Schnitzler an Gustav Schwarzkopf, 7. 8. 1926]{L04457 Arthur Schnitzler an Gustav Schwarzkopf, 7. 8. 1926}
\nopagebreak\mylabel{L04457v}
\rehead{ }\normalsize\beginnumbering\briefempfaengerindex{Schwarzkopf, Gustav@\textsc{Schwarzkopf, Gustav}!zzzSchnitzler, Arthur@\emph{von Arthur Schnitzler}!1926-08-071@{7. 8. 1926}|(be}
\toendnotes[C]{\smallbreak\pagebreak[2]}
\correspDesc{Versand  durch Arthur Schnitzler am 7. 8. 1926 in Adelboden
\newline{}Übermittlung  am 7. 8. 1926 in Adelboden
\newline{}Erhalt  durch Gustav Schwarzkopf im Zeitraum [8. 8. 1926 – 12. 8. 1926?] in Wien}\toendnotes[C]{\smallbreak}
\Standort{DLA, A:Schnitzler, HS.1985.1.1897.}
\physDesc{Bildpostkarte, 377 Zeichen
\newline{}Handschrift: Bleistift, lateinische Kurrent
\newline{}Versand: Stempel: »\nobreak{}\oindex{Adelboden@\textbf{Adelboden}|pwk}Adelboden (Berner Oberland), 7. VIII. 26, 17\nobreak{}«.  }\toendnotes[C]{\smallbreak}\pstart{}{\pb}\label{T_L04457-1v}\edtext{\textcolor{gray}{\textbf{A. S.}}}{\lemma{\textnormal{\emph{A. S.}}}\Cendnote{\textnormal{ovaler Absenderkleber}}}\label{T_L04457-1}\pend{}\pstart{}\textcolor{gray}{\textbf{WIEN, XVIII.}}\oindex{XVIII., Währing@\textbf{XVIII., Währing}|pw}\pend{}\pstart{}\textcolor{gray}{\textbf{STERNWARTESTR. 71}}\oindex{Wien@\textbf{Wien}!XVIII., Währing@\textbf{XVIII., Währing}!Sternwartestraße 71@\textbf{Sternwartestraße 71}|pw}\pend{}{\bigskip}\pstart{}{\pb}Austria\oindex{Österreich@\textbf{Österreich}|pw}\pend{}\pstart{}Hrn Guſtav Schwarzkopf\pend{}\pstart{}Wien I\oindex{I., Innere Stadt@\textbf{I., Innere Stadt}|pw}\pend{}\pstart{}Tiefer Graben 17\oindex{Wien@\textbf{Wien}!I., Innere Stadt@\textbf{I., Innere Stadt}!Tiefer Graben 17@\textbf{Tiefer Graben 17}|pw}\pend{}{\bigskip}
\pstart
           \noindent{}{\pb}\textcolor{gray}{\textbf{Nevada-Palace Hotel}}\oindex{Nevada Palace@\textbf{Nevada Palace}|pw},{\\}\textcolor{gray}{\textbf{Adelboden\oindex{Adelboden@\textbf{Adelboden}|pw}, Steghorn\oindex{Steghorn@\textbf{Steghorn}|pw} – Wildstrubel\oindex{Wildstrubel@\textbf{Wildstrubel}|pw}.}}\pend
           
\pstart
           \textcolor{gray}{\textbf{Phot. E. Gyger\pwindex{Gyger, Emanuel 6.\,3.\,1886 Achseten – 2.\,10.\,1951 Adelboden@\textsc{Gyger, Emanuel} (6.\,3.\,1886 Achseten – 2.\,10.\,1951 Adelboden)|pw}}}{ }{\\}\textcolor{gray}{\textbf{– Adelboden\oindex{Adelboden@\textbf{Adelboden}|pw} –}}\pend
           \vspace{1em}
\pstart
           \raggedleft{}{\pb}Adelboden\oindex{Adelboden@\textbf{Adelboden}|pw}, 7. 8. 926\pend
           \vspace{0.5em}
\pstart
           lieber Guſtav, leider hör ich nicht
      von Ihnen. Darf ich ein Wort
      nach \uline{Bern}\oindex{Bern@\textbf{Bern}|pw} postrest erhalten?
      (Oder telefoniren Sie an Frau
               \textcolor{gray}{Kleinbader}\pwindex{Kleinbader?? [Frau mit Telefonanschluss] @\textsc{Kleinbader?? [Frau mit Telefonanschluss]}|pw}.) – Hier könnte
      man sehr zufrieden sein, wenn{[}s{]}
      nicht gar zu viel regnete. \label{K_L04457-1v}\edtext{Dienstag}{\lemma{\textnormal{\emph{Dienstag}}}\Cendnote{\textnormal{Die Abreise erfolgte am Mittwoch, vgl. A. S.: \emph{Tagebuch}, 11. 8. 1926.}}}\label{K_L04457-1}
      fahr ich mit Heini\pwindex{Schnitzler, Heinrich 9.\,8.\,1902 Hinterbrühl – 12.\,7.\,1982 Wien@\textsc{Schnitzler, Heinrich} (9.\,8.\,1902 Hinterbrühl – 12.\,7.\,1982 Wien)|pw} fort, – über
      Bern\oindex{Bern@\textbf{Bern}|pw}, aber wohin weiſs ich noch 
      nicht. Schönste Grüße!\pend
           \pstart Ihr \spacefill\mbox{Arthur}.\pend{}\selectlanguage{ngerman}\endnumbering\briefempfaengerindex{Schwarzkopf, Gustav@\textsc{Schwarzkopf, Gustav}!zzzSchnitzler, Arthur@\emph{von Arthur Schnitzler}!1926-08-071@{7. 8. 1926}|)be}\mylabel{L04457h}
\begin{anhang}
\end{anhang}\newcommand{\dateiname}{L04457}\newcommand{\titel}{Arthur Schnitzler an Gustav Schwarzkopf, 7. 8. 1926}\newcommand{\editorInnen}{Herausgegeben von Jahnke, SelmaMüller, Martin Anton}\input{../tex-inputs/latex-leseansicht-abspann}
